\documentclass[aspectratio=169, 14pt]{beamer}
\usepackage{stampfly_slides}

\lesson{1}{モータ制御}{Motor Control}

\begin{document}

% ------------------------------------------------------------------
\begin{frame}
  \titlepage
\end{frame}

% ------------------------------------------------------------------
\begin{frame}{今日のゴール / Today's Goal}
  \begin{block}{目標}
    モータの番号と配置を理解し、PWM で個別制御する
  \end{block}

  \vspace{1em}

  \begin{itemize}
    \item モータ番号 M1--M4 の位置を覚える
    \item PWM デューティ比でモータ回転数を制御
    \item 安全な地上テストの方法を学ぶ
  \end{itemize}
\end{frame}

% ------------------------------------------------------------------
\begin{frame}{モータ配置 / Motor Layout}
  \begin{center}
    \resizebox{0.38\textwidth}{!}{\documentclass[border=10pt]{standalone}
\usepackage{tikz}
\usetikzlibrary{arrows.meta, positioning}

\begin{document}

\definecolor{cwcolor}{RGB}{0, 120, 200}    % Blue for CW
\definecolor{ccwcolor}{RGB}{220, 60, 20}   % Red for CCW
\definecolor{bodycolor}{RGB}{80, 80, 80}   % Dark gray
\begin{tikzpicture}[
    motor/.style={circle, draw, line width=1.5pt, minimum size=18mm, font=\large\bfseries},
    cw/.style={motor, draw=cwcolor, text=cwcolor},
    ccw/.style={motor, draw=ccwcolor, text=ccwcolor},
    label/.style={font=\small, text=bodycolor},
    arrow/.style={-{Stealth[length=3mm]}, line width=1.2pt},
  ]

  % Body (X-frame)
  \draw[bodycolor, line width=2pt] (-2.5, 2.5) -- (2.5, -2.5);
  \draw[bodycolor, line width=2pt] (2.5, 2.5) -- (-2.5, -2.5);

  % Center body
  \fill[bodycolor, rounded corners=3pt] (-0.6, -0.6) rectangle (0.6, 0.6);
  \node[font=\footnotesize, text=white] at (0, 0) {ESP32};

  % Front direction arrow
  \draw[arrow, bodycolor] (0, 1.0) -- (0, 2.0);
  \node[label, above] at (0, 2.0) {Front};

  % M1 - Front Right (CCW)
  \node[ccw] (m1) at (2.5, 2.5) {M1};
  \node[label, above right=-1mm] at (m1.north east) {FR};
  \draw[arrow, ccwcolor] ([shift={(-45:1.3)}]m1.center)
    arc[start angle=-45, end angle=225, radius=1.3cm];
  \node[font=\scriptsize, ccwcolor] at (4.2, 2.5) {CCW};

  % M2 - Rear Right (CW)
  \node[cw] (m2) at (2.5, -2.5) {M2};
  \node[label, below right=-1mm] at (m2.south east) {RR};
  \draw[arrow, cwcolor] ([shift={(45:1.3)}]m2.center)
    arc[start angle=45, end angle=-225, radius=1.3cm];
  \node[font=\scriptsize, cwcolor] at (4.2, -2.5) {CW};

  % M3 - Rear Left (CCW)
  \node[ccw] (m3) at (-2.5, -2.5) {M3};
  \node[label, below left=-1mm] at (m3.south west) {RL};
  \draw[arrow, ccwcolor] ([shift={(225:1.3)}]m3.center)
    arc[start angle=225, end angle=495, radius=1.3cm];
  \node[font=\scriptsize, ccwcolor] at (-4.2, -2.5) {CCW};

  % M4 - Front Left (CW)
  \node[cw] (m4) at (-2.5, 2.5) {M4};
  \node[label, above left=-1mm] at (m4.north west) {FL};
  \draw[arrow, cwcolor] ([shift={(135:1.3)}]m4.center)
    arc[start angle=135, end angle=-135, radius=1.3cm];
  \node[font=\scriptsize, cwcolor] at (-4.2, 2.5) {CW};

  % Legend
  \node[anchor=west, font=\small] at (-4.5, -4.5) {%
    \textcolor{cwcolor}{\rule{8pt}{8pt}} CW (clockwise) \quad
    \textcolor{ccwcolor}{\rule{8pt}{8pt}} CCW (counter-clockwise)};

\end{tikzpicture}
\end{document}
}
  \end{center}
  \begin{itemize}
    \item 対角のモータが同じ方向に回転（トルクバランス）
    \item M1(FR), M3(RL) = CW / M2(RR), M4(FL) = CCW
  \end{itemize}
\end{frame}

% ------------------------------------------------------------------
\begin{frame}{PWM とは / What is PWM?}
  \begin{columns}[T]
    \begin{column}{0.45\textwidth}
      \begin{center}
        \resizebox{!}{0.72\textheight}{\documentclass[border=10pt]{standalone}
\usepackage{tikz}
\usetikzlibrary{arrows.meta}

\definecolor{sigcolor}{RGB}{0, 120, 200}
\definecolor{dutycolor}{RGB}{0, 120, 200}
\definecolor{gridcolor}{RGB}{200, 200, 200}

\begin{document}
\begin{tikzpicture}[
    axis/.style={-{Stealth[length=2mm]}, line width=0.8pt},
    signal/.style={sigcolor, line width=1.5pt},
    fill area/.style={fill=dutycolor!15},
  ]

  \def\period{3}    % Period width in cm
  \def\height{1.2}  % Signal height in cm
  \def\yspacing{2.5} % Vertical spacing

  % --- 0% Duty ---
  \begin{scope}[yshift=3*\yspacing cm]
    \node[anchor=east, font=\bfseries] at (-0.3, \height/2) {0\%};
    \draw[axis] (0, 0) -- (2*\period+0.5, 0);
    \draw[axis] (0, 0) -- (0, \height+0.4);
    % Signal: always low
    \draw[signal] (0, 0) -- (2*\period, 0);
    % Period markers
    \foreach \i in {0, 1} {
      \draw[gridcolor, dashed] (\i*\period, 0) -- (\i*\period, \height);
    }
    \draw[gridcolor, dashed] (2*\period, 0) -- (2*\period, \height);
    \node[font=\scriptsize, below] at (\period/2, -0.2) {T};
    \draw[{Stealth[length=1.5mm]}-{Stealth[length=1.5mm]}, thin] (0, -0.15) -- (\period, -0.15);
  \end{scope}

  % --- 25% Duty ---
  \begin{scope}[yshift=2*\yspacing cm]
    \node[anchor=east, font=\bfseries] at (-0.3, \height/2) {25\%};
    \draw[axis] (0, 0) -- (2*\period+0.5, 0);
    \draw[axis] (0, 0) -- (0, \height+0.4);
    \foreach \i in {0, 1} {
      \pgfmathsetmacro{\xstart}{\i*\period}
      \pgfmathsetmacro{\xhigh}{\i*\period + 0.25*\period}
      \pgfmathsetmacro{\xend}{(\i+1)*\period}
      % Fill high portion
      \fill[fill area] (\xstart, 0) rectangle (\xhigh, \height);
      % Signal
      \draw[signal] (\xstart, \height) -- (\xhigh, \height)
                     (\xhigh, \height) -- (\xhigh, 0)
                     (\xhigh, 0) -- (\xend, 0);
      \ifnum\i=0
        \draw[signal] (\xstart, 0) -- (\xstart, \height);
      \fi
      % Grid
      \draw[gridcolor, dashed] (\xstart, 0) -- (\xstart, \height);
    }
    \draw[gridcolor, dashed] (2*\period, 0) -- (2*\period, \height);
    \draw[signal] (2*\period, 0) -- (2*\period, \height);
  \end{scope}

  % --- 50% Duty ---
  \begin{scope}[yshift=1*\yspacing cm]
    \node[anchor=east, font=\bfseries] at (-0.3, \height/2) {50\%};
    \draw[axis] (0, 0) -- (2*\period+0.5, 0);
    \draw[axis] (0, 0) -- (0, \height+0.4);
    \foreach \i in {0, 1} {
      \pgfmathsetmacro{\xstart}{\i*\period}
      \pgfmathsetmacro{\xhigh}{\i*\period + 0.5*\period}
      \pgfmathsetmacro{\xend}{(\i+1)*\period}
      \fill[fill area] (\xstart, 0) rectangle (\xhigh, \height);
      \draw[signal] (\xstart, \height) -- (\xhigh, \height)
                     (\xhigh, \height) -- (\xhigh, 0)
                     (\xhigh, 0) -- (\xend, 0);
      \ifnum\i=0
        \draw[signal] (\xstart, 0) -- (\xstart, \height);
      \fi
      \draw[gridcolor, dashed] (\xstart, 0) -- (\xstart, \height);
    }
    \draw[gridcolor, dashed] (2*\period, 0) -- (2*\period, \height);
    \draw[signal] (2*\period, 0) -- (2*\period, \height);
  \end{scope}

  % --- 100% Duty ---
  \begin{scope}[yshift=0cm]
    \node[anchor=east, font=\bfseries] at (-0.3, \height/2) {100\%};
    \draw[axis] (0, 0) -- (2*\period+0.5, 0);
    \draw[axis] (0, 0) -- (0, \height+0.4);
    % Fill entire area
    \fill[fill area] (0, 0) rectangle (2*\period, \height);
    % Signal: always high
    \draw[signal] (0, 0) -- (0, \height) -- (2*\period, \height);
    \foreach \i in {0, 1} {
      \draw[gridcolor, dashed] (\i*\period, 0) -- (\i*\period, \height);
    }
    \draw[gridcolor, dashed] (2*\period, 0) -- (2*\period, \height);
  \end{scope}

  % Title
  \node[font=\Large\bfseries, anchor=south] at (\period, 3*\yspacing + \height + 0.8)
    {PWM Duty Cycle};

  % Axis labels
  \node[font=\small, rotate=90, anchor=south] at (-1.5, 1.5*\yspacing)
    {Duty Cycle};

\end{tikzpicture}
\end{document}
}
      \end{center}
    \end{column}
    \begin{column}{0.5\textwidth}
      \vspace{1em}
      \begin{itemize}
        \item \textbf{Pulse Width Modulation}\\（パルス幅変調）
        \item \textbf{Duty 0\%} = 停止
        \item \textbf{Duty 100\%} = 最大回転
        \item API: \api{motor\_set\_duty(id, duty)}\\
              \small id=1--4, duty=0.0--1.0
      \end{itemize}
    \end{column}
  \end{columns}
\end{frame}

% ------------------------------------------------------------------
\begin{frame}{モータ制御 API}
  \begin{table}
    \centering
    \begin{tabular}{lll}
      \hline
      \textbf{関数} & \textbf{説明} & \textbf{引数} \\
      \hline
      \api{motor\_set\_duty(id, duty)} & 個別モータ設定 & id=1--4, duty=0.0--1.0 \\
      \api{motor\_set\_all(duty)} & 全モータ一括 & duty=0.0--1.0 \\
      \api{motor\_stop\_all()} & 全モータ停止 & --- \\
      \hline
    \end{tabular}
  \end{table}

  \vspace{1em}

  \begin{alertblock}{安全注意 / Safety}
    \warn{duty は 0.15（15\%）以下で！機体を手で押さえてテスト！}
  \end{alertblock}
\end{frame}

% ------------------------------------------------------------------
\begin{frame}[fragile]{実習: モータを順番に回す / Hands-on}
  \begin{lstlisting}[style=sfcpp, title={student.cpp}]
#include "workshop_api.hpp"
static uint32_t timer = 0;

void setup() {
    ws::print("Lesson 1: Motor Control");
}
void loop_400Hz(float dt) {
    timer++;
    // 800 ticks = 2 seconds at 400Hz
    int motor_id = (timer / 800) % 4 + 1;
    ws::motor_stop_all();
    ws::motor_set_duty(motor_id, 0.1f);
}
  \end{lstlisting}
\end{frame}

% ------------------------------------------------------------------
\begin{frame}{チェックポイント / Checkpoint}
  \begin{alertblock}{確認事項}
    \begin{itemize}
      \item[$\square$] M1 が FR（右前）のプロペラを回す
      \item[$\square$] M1$\to$M2$\to$M3$\to$M4 の順に回転する
      \item[$\square$] 各モータ 2 秒ごとに切り替わる
    \end{itemize}
  \end{alertblock}

  \vspace{1em}

  \begin{block}{次のレッスン}
    Lesson 2: コントローラ入力 --- スティックでモータを操作
  \end{block}
\end{frame}

\end{document}
